% Generated draft: requires mathematical and source review.
\documentclass[11pt]{article}
\usepackage[margin=25mm]{geometry}
\usepackage{fontspec}
\setmainfont{Latin Modern Roman}
\usepackage{amsmath,amssymb,mathtools,unicode-math}
\setmathfont{Latin Modern Math}
\usepackage{enumitem,needspace,xurl,hyperref}
\setlength{\parindent}{0pt}
\setlength{\parskip}{5pt}
\setlength{\emergencystretch}{3em}
\newcommand{\N}{\mathbb{N}}
\newcommand{\Z}{\mathbb{Z}}
\newcommand{\Q}{\mathbb{Q}}
\newcommand{\R}{\mathbb{R}}
\newcommand{\C}{\mathbb{C}}
\newcommand{\F}{\mathbb{F}}
\newcommand{\A}{\mathbb{A}}
\newcommand{\PP}{\mathbb{P}}
\newcommand{\E}{\mathbb{E}}
\newcommand{\eps}{\varepsilon}
\newcommand{\dd}{\,\mathrm{d}}
\begin{document}
\section*{UnsolvedMath ID: 2302009}
\textbf{Source problem code:} catalogue entry
\subsection{Definitions and mathematical statement}
The source asks for analogues of its Problems 2.6, 2.7 and 2.8 under the additional assumption that $f$ has another finite Picard value, with the explicit example
\[
 f(z)\ne0.
\]
Those three problems, their original function class, the earlier value referred to by ``another,'' and the formula labeled 2.1 are not included in the supplied record. Consequently the conclusions whose analogues are requested cannot be reconstructed.

For terminology only, if $f$ is a nonconstant meromorphic function on $\C$, a finite omitted value is a number $b\in\C$ such that $f(z)\ne b$ for every $z$; some Picard-exceptional terminology instead permits finitely many preimages. The source's displayed zero-free example is unambiguously an omitted-value condition, but its intended general Picard-value convention is unspecified. In the same plane-meromorphic setting,
\[
 \delta(\infty,f)=1-\limsup_{r\to\infty}\frac{N(r,\infty;f)}{T(r,f)},
\]
where $N$ counts poles with multiplicity and $T$ is the Nevanlinna characteristic. Thus deficiency one means the displayed counting-to-characteristic ratio tends to zero in limsup, not necessarily that there are no poles. Finite versus infinite order can be expressed by whether $\limsup_{r\to\infty}\log^+T(r,f)/\log r$ is finite or infinite. These definitions explain the quoted conditions without silently asserting that the omitted original problems use precisely this domain.

The source reports that, with the extra Picard-value condition and $\delta(\infty,f)=1$, its formula 2.1 remains true in finite order, citing Edrei and Fuchs, but need not hold in infinite order, citing Gol'dberg and Ostrovskii. The formula itself is absent. No replacement formula or complete theorem is guessed from this report.
\subsection{Short English statement}
Study the source's three preceding questions when an additional finite value is omitted or Picard-exceptional. The supplied record does not contain those questions or their referenced formula, so the actual requested analogues remain undefined. It distinguishes a reported finite-order result from a failure of the cited formula in general infinite order. The supplied triage identifies the infinite-order case as open in Hayman's 2018 update and reports no recent resolution located; this is an attributed dated assessment, not a new status review.
\subsection{Sources}
\begin{itemize}
\item UnsolvedMath permanent ID 2302009, source code catalogue entry, ``Research Problems in Function Theory --- Problem 2.9.'' Registry curation and metadata: Ulam AI and UnsolvedMath Contributors, CC BY 4.0; updated 2026-08-06.
\item Hayman, \emph{Research Problems in Function Theory} (2018 source update), Problem 2.9, referring internally to Problems 2.6--2.8 and formula 2.1: \url{https://arxiv.org/abs/1809.07200}. Supplied source-TeX extraction accessed 2026-07-29. The metadata retains the 2018 update for status triage and marks status/release and rights review NEEDS\_REVIEW.
\item Edrei and Fuchs, and Gol'dberg and Ostrovskii, as named in the source statement. Complete references and the corrupted parenthetical reference contents are not supplied; none are reconstructed by inference.
\end{itemize}
\end{document}
